\usepackage[a4paper, margin=2cm]{geometry} % Formatage de la page
\usepackage[english,french]{babel} % Langues, dont la principale en dernière
\usepackage{graphicx} % Images
\usepackage{placeins} % Float Barriers
\usepackage{xcolor} % Couleurs
\usepackage[most]{tcolorbox} % Encadrés personnalisables
\usepackage{lstautogobble} % Gestion des indentations pour le code dans les encadrés
\usepackage{tikz} % Formes
\usepackage{tikzpagenodes} % Points d'ancrage supplémentaires pour TikZ
\usepackage{amssymb} % Signes supplémentaires (dont checkmark)
\usepackage{parskip} % Présentation plus aérée
\usepackage{enumitem} % Présentation des listes
\usepackage[autostyle]{csquotes} % Gestion des citations
\usepackage[backend=biber]{biblatex} % Bibliographie
\usepackage[colorlinks=true, linkcolor=blue, urlcolor=blue, citecolor=blue]{hyperref} % Liens hypertextes

\makeatletter
\g@addto@macro\UrlBreaks{\do\-} % Permet au \path de se casser au niveau des tirets
\makeatother